\chapter{Physics Is Not Data Science}
\label{physics-is-not-data-science}
\epigraph{Not being able to define a thing precisely is not the same failure 
as not understanding it.}{ R. Feynman}
 Physics survives without precise definitions because it replaces them with
 measurement procedures. This chapter argues that measurement procedures are
 the \emph{floor} of physics, not its ceiling --- and that the gap between the
 floor and the ceiling is exactly what separates a physicist from a
 curve-fitting algorithm.

The claim in the epigraph above rests on two sentences Feynman actually said,
in different lectures, about different things -- and they are routinely
collapsed into one. They should not be.

``We cannot define anything precisely'' was said about ordinary words like
\emph{motion} -- a confession that physics does not wait for philosophical
certainty about what a thing \emph{is} before it starts measuring it precisely.
That is an epistemic strategy: start with measurement, because measurement is
checkable and definition-by-essence is not.

``I think I can safely say that nobody understands quantum mechanics'' was said
in a completely different spirit.  Feynman had just finished pointing out that
although newspapers once claimed only twelve people understood relativity, that
was never really true --- once Einstein's paper was published, plenty of people
came to understand it, in the ordinary sense of being able to picture what it
says. Quantum mechanics, he then says, is different: nobody has ever been able
to reduce it to a picture built out of everyday experience, the way you can
eventually picture curved spacetime with enough effort. That is not a statement
that physicists cannot use quantum mechanics, predict with it, or reason about
it --- they do all three, correctly, every day, and Feynman himself spent the
rest of that same lecture doing exactly that. It is a much narrower claim: no
one has found a \emph{classical, intuitive mechanism} underneath it.
Understanding, in that specific and demanding sense, is missing. Understanding
in the working sense --- knowing the structure, deriving its consequences,
trusting it enough to build a transistor or a laser on top of it --- is not.

This chapter is about the difference between those two kinds of understanding,
and about a claim worth taking seriously: that pure operationalism, on its own,
cannot tell them apart. If all that mattered were successful measurement and
prediction, there would be no principled way to say that a physicist
understands gravity more deeply than a weather-prediction algorithm understands
rain. Physics needs, and has, a further standard.

Return to  account of measurement. An operational definition says: a quantity
is whatever a specified procedure outputs. This is enough to make a concept
scientific --- checkable, shareable, free of unfalsifiable metaphysics. It is
not enough, on its own, to distinguish two very different kinds of success.
\begin{enumerate}
\item A theory that predicts a measurement because it encodes something
    structural about how the world is put together: a small number of
        principles from which the measured value can be \emph{derived}.
\item A model that predicts a measurement because it has been fitted, with
    enough free parameters, to enough past instances of that same measurement,
        with no claim about structure at all.
\end{enumerate}

Both can satisfy the operationalist's only stated requirement: the output
matches what the ruler, clock, or detector reads. Bridgman's operationalism,
taken as a complete philosophy of physics rather than as a starting discipline,
has no vocabulary for preferring one over the other. Yet every working
physicist does prefer one over the other, strongly and immediately. The rest of
this chapter tries to say, precisely, why.

The clearest demonstration Feynman ever gave of this distinction was not a
slogan; it was a worked example.He points out that Newton's law of gravitation
can be written in at least three completely different-looking forms.

\begin{description}
\item[Action at a distance.] Each mass reaches out across empty space and pulls
    on every other mass directly:
\[
  F = \frac{G m_1 m_2}{r^2}.
\]
\item[The field method.] Every mass creates a field filling all of space; a
    second mass simply responds to the local field value at its own location,
        with no need to reach anywhere:
\[
  \nabla^2 \phi = 4\pi G \rho, \qquad \vec F = -m\,\nabla \phi.
\]
\item[The principle of least action.] Of all the conceivable paths a particle
    could take between two points, nature selects the one that makes a
        particular quantity, the action $S$, stationary:
\[
  \delta S = \delta\!\int (T - V)\, dt = 0.
\]
\end{description}

For every experiment classical mechanics can perform, these three formulations
issue identical predictions. Bridgman's yardstick sees no difference between
them at all: same measurement in, same number out, for all three, always. And
yet, as Feynman insists, they are not the same idea wearing different clothes.
They organize your intuition differently, suggest different questions, and ---
this is the decisive part --- they are not equally good starting points for the
\emph{next} theory.

When physics moved from classical to quantum mechanics, the
action-at-a-distance picture and the field-at-a-point picture both had to be
abandoned or radically reworked. The least-action formulation, by contrast,
generalized almost gracefully: Feynman's own path-integral formulation of
quantum mechanics, developed in 1948, is built by summing a phase
$e^{iS/\hbar}$ over every conceivable path, weighting the classical path of
least action as the one that survives in the limit $\hbar \to 0$. Three
formulations, operationally indistinguishable in 1687 and still operationally
indistinguishable today for any classical experiment --- but only one of them
was carrying the structure that quantum mechanics needed. Operational
equivalence in the old domain said nothing at all about generative power in the
new one.

That gap is what \emph{understanding} means in physics. It is not merely
getting today's number right. It is possessing a formulation with enough
internal structure that it tells you, in a principled way, what to try next.

\begin{figure}[ht]
\centering
\begin{tikzpicture}[every node/.style={font=\footnotesize}, >=stealth]
  \node[draw, rounded corners,  minimum width=2.7cm, minimum height=0.85cm, align=center] (center)
    {Same predictions:\\ classical gravitation};

  \node[draw, rounded corners, minimum width=2.1cm, minimum height=0.75cm, align=center, above left=0.9cm and 0.1cm of center] (force) {Action at\\ a distance};
  \node[draw, rounded corners, minimum width=2.1cm, minimum height=0.75cm, align=center, above right=0.9cm and 0.1cm of center] (field) {Field\\ method};
  \node[draw, rounded corners, minimum width=2.1cm, minimum height=0.75cm, align=center, below=1.0cm of center] (action) {Principle of\\ least action};

  \draw[<->, thick] (force) -- (center);
  \draw[<->, thick] (field) -- (center);
  \draw[<->, thick] (action) -- (center);

  \node[draw, dashed, rounded corners, minimum width=2.1cm, minimum height=0.65cm, align=center, left=0.7cm of force, yshift=-0.25cm] (deadforce) {no natural\\ generalization};
  \node[draw, dashed, rounded corners, minimum width=2.1cm, minimum height=0.65cm, align=center, right=0.7cm of field, yshift=-0.25cm] (deadfield) {no natural\\ generalization};
  \node[draw, rounded corners,  minimum width=2.9cm, minimum height=0.85cm, align=center, below=1.0cm of action] (qm) {Quantum mechanics\\ via path integrals (1948)};

  \draw[->, thick] (force) -- (deadforce);
  \draw[->, thick] (field) -- (deadfield);
  \draw[->, thick] (action) -- (qm);
\end{tikzpicture}
\caption{Three formulations of Newton's law of gravitation, operationally
    identical for every classical measurement --- yet only the least-action
    formulation carried the structure needed to generalize to quantum
    mechanics. Operational equivalence in one domain does not imply equal
    understanding, or equal reach.}
\label{fig:three-ways}
\end{figure}

The philosopher and computer scientist Judea Pearl has proposed a useful
vocabulary for the same distinction, developed for a very different purpose:
distinguishing statistical, correlational reasoning from causal reasoning. He
describes three rungs of a ladder.

\begin{center}
\begin{tabular}{@{}>{\raggedright\arraybackslash}p{2.6cm} >{\raggedright\arraybackslash}p{4.6cm} >{\raggedright\arraybackslash}p{4.6cm}@{}}
\textbf{Rung} & \textbf{Typical question} & \textbf{What answers it} \\
Association & What does a value of $X$ tell me about $Y$? & A fit to past data; correlation \\
Intervention & What happens to $Y$ if I \emph{change} $X$? & A causal / mechanistic model \\
Counterfactual & What would $Y$ have been, had $X$ been different? & A model
    that supports reasoning beyond any data actually collected \\
\end{tabular}
\end{center}

A purely operational, purely data-fitted model can climb no higher than the
first rung by construction: it was built to reproduce association between
measurement and measurement, and nothing in its construction licenses a claim
about what would happen under conditions it has never seen. A physical theory
with real structure --- Newton's law derived from a principle, not merely
fitted to Tycho Brahe's tables --- operates on the second and third rungs
almost as a matter of course. It tells you what happens if you double the mass
of the Sun, a question no data set has ever answered, because the answer is
derivable rather than looked up.

History supplies the record of this pay-off. The existence and orbit of Neptune
were predicted, in 1846, from irregularities in Uranus's orbit, before anyone
pointed a telescope at the right patch of sky. The positron was predicted by
Dirac's equation in 1928, purely from requiring the equation to be consistent,
four years before one was observed in a cloud chamber. Gravitational waves were
predicted in 1916 and detected almost exactly a century later, in 2015, with
the detected waveform matching the theoretical prediction for two merging black
holes to remarkable precision. None of these were interpolations inside a
training set. Each was an extrapolation, licensed by structure, into territory
no measurement had ever visited --- exactly the kind of claim Pearl's first
rung cannot support, and exactly what physicists mean when they say they
\emph{understand} gravity, or the electron, rather than merely being able to
predict its next measured value.

\section{Physics and Data Science}

This is the precise place where physics and data science part ways, and it is
worth stating the comparison without exaggerating it in either direction.

Train a sufficiently flexible statistical model on centuries of recorded
planetary positions and it will predict tomorrow's positions beautifully ---
quite possibly to a precision rivaling Newton's own equations, for exactly the
regime it was trained on. It has, in the fullest operational sense, learned to
measure and predict planetary motion. What it has not done is acquire a reason
to trust its own output the moment the situation changes: a new comet on a
hyperbolic orbit, a spacecraft under active thrust, a binary star system where
the gravitational field is far stronger than anything in the training data. A
model built on correlation has no internal signal that it has left its depth
--- it will output a confident number regardless. Newton's law, by contrast,
comes with a built-in, structural account of exactly where and why it should be
expected to fail: at speeds approaching $c$, or fields strong enough that
spacetime curvature can no longer be neglected. Knowing your own boundary of
validity is itself a form of understanding, and it is precisely what falls out
of having a mechanism rather than a fit.

None of this makes data-driven methods unscientific, and it would be a mistake
to read this chapter that way. Statistical inference is indispensable precisely
where no mechanistic theory yet exists, and some of the most interesting
current work in physics itself uses machine learning as a tool for discovery
--- to search for patterns in particle-collider data, or in gravitational-wave
signals, that a human physicist might not think to look for. Tellingly, though,
the most physically successful of these tools are increasingly built to encode
known structure directly --- conservation laws, symmetries, the least-action
principle itself --- rather than learning everything from scratch. That design
choice is itself an admission of this chapter's argument: bare correlation
extrapolates badly, and the fix physicists reach for is always the same one
Feynman reached for in 1948 --- find the formulation that carries structure,
not just the one that fits the numbers.


Chapter~\ref{reality-and-models} ended with a motto built from three of
Feynman's remarks: physics
does not need to know what space, time, or motion \emph{are}, only how to
measure them. That motto is correct as far as it goes, and it is where physics
has to start. This chapter has argued that it is not where physics is content
to stop.

A physicist who has only an operational definition and an accurate prediction
has done real, indispensable work --- exactly the work 
Chapter~\ref{reality-and-models} celebrated. A
physicist who additionally has a structure --- a small set of principles from
which that prediction, and a thousand others never yet measured, can all be
derived --- has done something further, and it is that further thing which
earns the word \emph{understand}. The motto of this chapter is the necessary
second half of the first:

\emph{We measure a thing before we understand it. We understand it once our
account of it lets us predict what we have never measured at all.}

\section{The Physicist as Intuitive Bayesian}

A physicist does not confront data from a position of ignorance. Before a
single measurement is taken, she has already restricted attention to a family
of models $\{M_\theta : \theta \in \Theta\}$ --- Newton's law with an unknown
spring constant, a resonance of unknown mass and width, a cosmological model
with unknown $\Omega_\Lambda$. This restriction is not neutral. It is an act of
assigning prior plausibility zero to the overwhelming majority of functions
that could in principle be fit to any given data set, and non-negligible
plausibility to a small, structured subset selected by symmetry,
renormalizability, dimensional analysis, or simply what has worked before. This
is the content of a Bayesian prior, whether or not it is ever written as
$\pi(\theta)$. The thesis of this section is that physics, as actually
practiced, is Bayesian inference over a self-imposed model space --- even when,
and especially when, the explicit statistical vocabulary used to report the
result is frequentist.

\subsection{Model space as an implicit prior}

\begin{definition}
A \emph{model space} $\mathcal{M}$ is the set of hypotheses a physicist is
    willing to entertain prior to seeing the data at hand, together with a
    measure $\pi$ of relative plausibility over that set --- possibly only
    ordinal, and rarely written explicitly.
\end{definition}

Formally, inference within a chosen model proceeds by
\begin{equation}
P(\theta \mid D, M) \;=\; \frac{P(D \mid \theta, M)\,\pi(\theta \mid M)}
    {\displaystyle\int_{\Theta} P(D \mid \theta', M)\,\pi(\theta' \mid M)\, d\theta'}
    \,, \qquad \theta \in \Theta_M,
\end{equation}
but the physicist's real act of judgment happens one level up, in the choice of
$M$ itself from the space of all conceivable models $\mathcal{M}$. Requiring
Lorentz invariance, gauge invariance, renormalizability, or simply a low-order
polynomial in the fields are all instances of setting $\pi(M') = 0$ for models
$M'$ outside the chosen family --- an extremely informative prior, imposed
before any likelihood is ever evaluated, and rarely defended in the
probabilistic language it deserves.

\begin{remark}
Occam's razor, as actually practiced in theoretical physics --- preferring the
    model with fewer free parameters, absent a specific reason to prefer
    otherwise --- is not a separate methodological principle standing outside
    probability theory. It is a description of a prior that penalizes model
    complexity, and it has a standard formal counterpart in Bayesian model
    comparison: the automatic Occam penalty carried by the marginal likelihood
\begin{equation}
P(D \mid M) \;=\; \int_{\Theta} P(D \mid \theta, M)\,\pi(\theta \mid M)\, d\theta,
\end{equation}
which is suppressed for models whose parameter space wastes prior mass on
    regions the data rule out.
\end{remark}

\subsection{Sequential updating across experiments}

The history of confidence in General Relativity is not well described as a
sequence of independent frequentist tests, each beginning from total ignorance
and returning a verdict discarded once the next experiment starts. Each
successive test --- Mercury's perihelion, light bending, gravitational
redshift, binary-pulsar timing, the direct detection of gravitational waves ---
is treated as tightening an already substantial body of belief, and a single
anomalous result is weighed \emph{against}, not substituted for, everything
already established. This is posterior-to-prior updating,
\begin{equation}
\pi_n(\theta) \;\propto\; P(D_n \mid \theta)\,\pi_{n-1}(\theta),
\end{equation}
carried out informally, paper by paper, over decades, whether or not any
individual paper ever computes a posterior.

\subsection{Confidence intervals: frequentist construction, Bayesian use}

A frequentist $100(1-\alpha)\%$ confidence interval $I(D)$ is defined by the
coverage property
\begin{equation}
P_\theta\big(\theta \in I(D)\big) = 1-\alpha \qquad \text{for every fixed } \theta,
\end{equation}
a statement about the behavior of the \emph{procedure} across a hypothetical
ensemble of repetitions of the experiment at the same true $\theta$. By
construction it makes no probabilistic statement about whether the true value
lies in the one particular interval actually obtained from the one experiment
that was run.

\begin{remark}
This is almost never how such an interval is read once quoted. ``The top-quark
    mass is $172.5 \pm 0.7\ \mathrm{GeV}$'' is universally understood,
    including by the physicists reporting it, as licensing something close to
    ``there is roughly a 68\% chance the true mass lies in this range'' --- a
    statement about degree of belief in a fixed but unknown constant,
    conditioned on the one data set actually obtained. That is a Bayesian
    credible interval, not the frequentist object formally computed. The two
    coincide numerically under a flat prior in a wide class of problems, which
    is precisely why the conflation goes unnoticed.
\end{remark}

\subsection{Systematic uncertainty as a nuisance-parameter prior}

A detector's calibration constant, an integrated luminosity, a background
normalization --- these are not resampled across hypothetical repetitions of
the universe; they are fixed, unknown quantities specific to the one apparatus
and the one run at hand. Assigning them an error distribution and
marginalizing,
\begin{equation}
P(\theta \mid D) \;=\; \int P(\theta \mid D, \lambda)\,\pi(\lambda)\, d\lambda,
\end{equation}
over a nuisance parameter $\lambda$ --- calibration, background rate,
theoretical cross-subsection uncertainty --- is exactly Bayesian
nuisance-parameter marginalization. There is no frequentist repeated-sampling
story under which ``the calibration constant of this detector, on this day, has
a 90\% confidence interval'' is a meaningful statement: there is only one
calibration constant, and one day. Experimental physicists carry out this
marginalization routinely, correctly, and almost never under that name.

\subsection{Case study: the $5\sigma$ threshold and the Higgs boson}

The 2012 discovery announcement was couched entirely in frequentist terms: a
local significance in excess of $5\sigma$, i.e.\ a $p$-value under the
background-only null hypothesis below roughly $3\times10^{-7}$. This is, by
construction, \emph{not} a posterior probability that the Higgs boson exists,
and not a statement about $P(\text{Higgs} \mid D)$ at all --- a point
statisticians raised at the time and which physicists mostly reasoned past
rather than through.

\begin{remark}
What made the announcement compelling was never the $p$-value in isolation. It
    was the $p$-value combined with decades of prior theoretical commitment ---
    electroweak precision fits, unitarity arguments against a Higgsless
    Standard Model, a mass window already narrowed by indirect constraints ---
    that gave the community a strong prior concentrated near the observed mass
    before ATLAS and CMS reported anything. The frequentist significance
    functioned, in practice, as a likelihood-ratio-like ingredient fed into an
    unstated Bayesian update, not as a free-standing verdict evaluated against
    a flat prior over ``a new particle exists somewhere in an otherwise
    unconstrained parameter space.'' A $5\sigma$ excess with no such prior
    theoretical support is treated far more skeptically by the same community,
    appropriately --- itself evidence that the underlying reasoning is Bayesian
    even while the reported statistic is not.
\end{remark}

\subsection{Real-world content versus ideal-world ensembles}

This section's thesis connects directly to the real-world/ideal-world
distinction developed earlier in this book. The frequentist confidence interval
is defined via an ensemble of hypothetical repetitions of the experiment that
were never run, and, for a one-off cosmological measurement or a unique
collider run at a given energy, could not in principle be run again
identically. That ensemble is precisely the kind of ideal-world construct
examined in the chapter on numerical calculation: a formally well-defined
object, useful for calibrating a procedure, that does not itself denote
anything that has happened in the real world. The Bayesian posterior, by
contrast, is a degree of belief about the one already-realized data set and the
one fixed constant of nature the physicist actually cares about --- real-world
content, in this book's sense, with no further translation step required to
make contact with it.

\begin{corollary}
A physicist who reports a frequentist interval but reasons, updates confidence,
    and makes decisions as though it were a posterior is not committing an
    error of arithmetic; she is correctly tracking real-world content while
    mislabeling the ideal-world formalism used to obtain it --- the same
    species of world-mixing diagnosed, from the opposite direction, in the
    chapter on Euclidean lattice calculation.
\end{corollary}

\subsection{Closing}

None of this is an argument that frequentist tools are dispensable. Profile
likelihoods, $p$-values, and coverage guarantees remain the disciplined
machinery by which a model survives contact with data without the analyst's
prior smuggling in the desired answer unexamined. The claim is narrower, and,
the author hopes, more interesting: the physicist's actual epistemic practice
--- restricting the hypothesis space before data arrives, updating confidence
across a lifetime of experiments rather than resetting it with each one,
treating fixed unknowns as objects that can bear probability --- is Bayesian in
structure, regardless of which formal statistical vocabulary is used to report
the result.

